\documentclass{article}

\usepackage[spanish, es-tabla]{babel}
\usepackage{graphicx}
\usepackage[hidelinks]{hyperref}
\usepackage{amssymb}
\usepackage{amsmath}

\spanishdecimal{.}

\begin{document}
\begin{titlepage}
{\includegraphics[width=1\textwidth]{./img/logoFC.png}\par}
\vspace{1cm}
\centering
{\bfseries\huge Universidad Nacional Autónoma de México \par}
\vspace{1cm}
{\scshape\huge Facultad de Ciencias \par}
\vspace{2cm}
{\scshape\Huge Tarea 2 - Reporte\par}
\vspace{2cm}
{\itshape\LARGE Computo Evólutivo \par}
\vfill
{\large Autores: \par}
{\large Diego García V. \par}
{\large Julio Vázquez A. \par}
\vfill
{\Large \today \par}
\end{titlepage}

\section*{Problema de la Mochila 0-1 }
Dada una colección de \textit{n} objetos y una mochila, tal que:

\begin{center}
$p_j = $ beneficio (ganancia) del objeto $j$

$w_j = $ peso del objeto $j$

$c = $ capacidad de la mochila
\end{center}

sujeto a:

\begin{equation*}
\sum_{j=1}^n w_{j}x_{j} \leq c
\end{equation*}

donde:


\begin{equation*}
 x_{j}=
\begin{cases}
1, & \text{si el elemento \textit{j} es seleccionado }\\
0,& \text{e.o.c }
\end{cases}
\end{equation*}




\subsection*{Descripción de la implementación}
%TODO

\subsection*{Descripción del esquema de representación de soluciones}
%TODO

\subsubsection*{Ventajas}
%TODO

\subsubsection*{Desventajas}
%TODO

\subsection*{Descripción de la función de evaluación}
%TODO

\subsubsection*{Función de evaluación de soluciones}
%TODO

\subsection*{Descripción del generador de soluciones aleatorias}
%TODO

\subsubsection*{Representación de soluciones}
%TODO

\subsection*{Descripción de la función u operador de vecindad}
%TODO

\subsection*{Recocido simulado}
%TODO

\subsubsection*{Descripción del esquema de enfriamiento}
%TODO

\subsection*{Tablas}
%TODO

\subsubsection*{Resultados obtenidos}
%TODO

\newpage
%TODO
\section*{Referencias}
\begin{enumerate}
	\item 
\end{enumerate}

\end{document}
